\documentclass[12pt]{article}
\usepackage[utf8]{inputenc}
\usepackage[T1]{fontenc}
\usepackage{hyperref}
\title{Malaysian Sign Language Language Resources for K-12 Curriculum Development}
\author{Compiled by Claude Sonnet 5.5 (claude-sonnet-5-5) for Andr\'as Kornai}
\date{}
\begin{document}
\maketitle

\section{Introduction}
Malaysian Sign Language (Bahasa Isyarat Malaysia, BIM; ISO 639-3 xml; Glottocode mala1412) is the sign language of the Malaysian deaf community. The search results describe it as a full visual language, formally recognised in 2008 under the Persons with Disabilities Act, with roots in regional sign varieties and some influence from American Sign Language. The same sources say Kod Tangan Bahasa Malaysia (KTBM), a manually coded form of Malay, is the system the Ministry of Education recognises for teaching Malay to deaf students. That claim comes from a secondary source (Wikipedia). My searches found little in the way of ready-made K-12 materials in BIM. This report lists only what the searches returned.

\section{Primary Educational Resources}
No BIM primer, textbook or school curriculum turned up in the searches. The closest educational items are two research papers from the University of Malaya journal \emph{JuKu} on using BIM to teach basic Islamic obligations (PAFA) to deaf learners. They report that BIM helped learners understand the content, and they identify a lack of signs for Islamic terms and a lack of a BIM curriculum for PAFA.
\begin{itemize}
\item Importance of BIM in teaching PAFA to the deaf community: \url{https://ejournal.um.edu.my/index.php/JUKU/article/view/8136}
\item Content design of a BIM curriculum for PAFA: \url{https://ejournal.um.edu.my/index.php/JUKU/article/view/8165}
\end{itemize}

\section{Linguistic and Reference Materials}
A search-result summary mentions a sign book published by the Malaysian Federation of the Deaf, but I could not confirm its availability and found no URL for it. No current BIM dictionary or sign database turned up.
\begin{itemize}
\item Overview of the language (secondary source): \url{https://wikipedia.com/wiki/ISO_639:xml}
\end{itemize}

\section{Religious and Community Resources}
\begin{itemize}
\item Taylor's College article on BIM and the deaf community in Malaysia: \url{https://college.taylors.edu.my/en/student-life/news/2025/bim-breaking-the-silence-in-malaysia.html}
\end{itemize}
The PAFA studies listed above also bear on religious education.

\section{Digital and Audio Resources}
No video corpora, apps or sign recordings turned up in the searches.

\section{Curriculum Development Context}
The evidence found is thin. A curriculum in BIM would need community-led development with Malaysian deaf organisations. The Malaysian Federation of the Deaf (Persekutuan Orang Pekak Malaysia) and the Ministry of Education's special education division are natural first contacts, but I did not find their pages in these searches. The documented gaps are signs for specialised vocabulary and curriculum design in BIM. The KTBM policy suggests that formal schooling currently favours a manually coded form of Malay over BIM.

\end{document}